\documentclass[11pt]{article}
\usepackage[margin=1in]{geometry}
\usepackage{amsmath,amssymb,amsthm}
\newtheorem{lemma}{Lemma}
\newcommand{\N}{\mathcal N}
\newcommand{\rank}{\operatorname{rank}}
\title{Independent audit of the uniform block-rank argument}
\date{November 2026}
\begin{document}
\maketitle

Fix $0<\alpha<1$ and put
\[
 K(s)=\frac{\sinh(\alpha s/2)}{\sinh(s/2)},\qquad
 a=\frac{1-\alpha}{2},\quad b=\frac{1+\alpha}{2}.
\]
The argument below checks the proposed finite-measure rank lemma and
its summation. It does not use convergence of finite graphs.

\begin{lemma}
Let $\mu$ be any positive finite measure of mass $m$ supported on an
interval of length $h\ge1$. For every $R\ge1$, the integral operator
with kernel $K(x-y)$ on $L^2(\mu)$ has an approximant of rank at most
\[
 \left(\frac{h+1}{2\pi^2}+4\right)R+O_\alpha(h+1)
\]
and operator-norm error at most $m e^{-R}$.
\end{lemma}
\begin{proof}
Write $T=h+1\ge2$ and periodize $K$:
\[
 P_T(s)=\sum_{r\in\mathbb Z}K(s+rT).
\]
The Fourier symbol
\[
 \widehat K(\xi)=
 \frac{2\pi\sin(\pi\alpha)}{\cosh(2\pi\xi)+\cos(\pi\alpha)}
\]
is bounded by $C_\alpha e^{-2\pi|\xi|}$. Poisson summation, or direct
integration of the absolutely convergent periodization, gives Fourier
coefficients $T^{-1}\widehat K(2\pi k/T)$. Truncating to $|k|\le q$
therefore gives rank at most $2q+1$ and uniform kernel error at most
\[
 \frac{2C_\alpha}{T}
 \frac{e^{-4\pi^2(q+1)/T}}{1-e^{-4\pi^2/T}}
 \le C'_\alpha e^{-4\pi^2q/T}.
\]
Here the prefactor is uniformly bounded for $T\ge2$.

For $t>0$ the exact absolutely convergent expansion is
\[
 K(t)=\sum_{j\ge0}\bigl(e^{-(j+a)t}-e^{-(j+b)t}\bigr).
\]
Thus, for $|s|\le h$,
\[
 P_T(s)-K(s)=\sum_{j\ge0}\left[
 d_{j+a}(e^{(j+a)s}+e^{-(j+a)s})
 -d_{j+b}(e^{(j+b)s}+e^{-(j+b)s})\right],\quad
 d_c=\frac{e^{-cT}}{1-e^{-cT}}.
\]
Each retained value of $j$ contributes at most four rank-one kernels
after setting $s=x-y$. The tail beginning at $j=J$ has absolute value
at most
\[
 \sum_{j\ge J}\sum_{c\in\{j+a,j+b\}}
 \frac{2e^{-c(T-h)}}{1-e^{-cT}}
 \le C''_\alpha e^{-J},
\]
because $T-h=1$ and $T\ge2$. Taking
\[
 q=\left\lceil\frac{T(R+c_\alpha)}{4\pi^2}\right\rceil,
 \qquad J=\lceil R+c_\alpha\rceil
\]
with fixed sufficiently large $c_\alpha$ gives kernel error at most
$e^{-R}$ and the displayed rank bound. A kernel bounded in absolute
value by $e^{-R}$ has Hilbert--Schmidt norm at most $m e^{-R}$,
including for atomic measures. This proves the operator-norm statement.
\end{proof}

For completeness, suppose a measure on $[0,\infty)$ satisfies
$\mu([s,\infty))\le A e^{-s/2}$ and has bounded total mass. Put
$L=\log(1/\varepsilon)$, $h=\sqrt L$, $S=2L+8\log L$, and
$M=\lceil S/h\rceil$. Partition into $M$ intervals of length $h$
and a tail. The corresponding positive operator is bounded above by
$(M+1)$ times its block diagonal pinching. The tail block has norm
at most $\alpha A e^{-Mh/2}\le\alpha A\varepsilon L^{-4}$;
it therefore contributes no count at threshold
$\varepsilon/(M+1)$ for large $L$.

The $i$th finite block has mass at most $A e^{-ih/2}$. Apply the lemma
with
\[
 R_i=\max\{1,L-ih/2+\log[A(M+1)]\}.
\]
Its error is at most $\varepsilon/(M+1)$. Uniformly as $L\to\infty$,
\[
 \sum_{i=0}^{M-1}R_i
 \le\frac{L^2}{h}
 +O_A\left(L+\frac{L\log(M+1)}h+M\right).
\]
For example this follows by comparing the decreasing positive part
of $L+O_A(\log M)-ih/2$ with its triangle integral. Hence
\[
 \N(\varepsilon)
 \le \frac{L^2}{2\pi^2}+O_{\alpha,A}(L^{3/2}+L\log L).
\]
All constants are independent of the measure and of its number of
atoms. The same argument applies to a discrete graph leverage measure
after $s=\log(4/\lambda)$, for which $A=2$ is supplied by
$\nu((0,\delta])\le\sqrt\delta$.

\end{document}
